
\section{Preliminaries}
\label{sec:pre}
\heading{Notations.}
In this paper, algorithms are modeled as \emph{functions} or \emph{(streaming) word-RAMs}. A function represents a circuit with a specified domain and range, and we measure its complexity by its depth (i.e., the parallel running time). A word-RAM has access to memory and registers, each holding one word. Its running time is measured in the number of basic operations (e.g., addition, subtraction, multiplication, bit-shifting, and memory access). A word-RAM is a streaming one if it receives inputs in a streaming manner and its local memory is measured in bits.

%Rewinding random bits used by RAMs is not permitted, so if an algorithm wants to access previously used random bits it must store them.

We write $a \getsr \A(b)$ (respectively, $a=\A(b)$) to denote the random variable output by a probabilistic (respectively, deterministic) algorithm $\A$ on input $b$.
By $x \getsr \BS$ we denote the process of sampling an element $x$ from a set or distribution $\BS$ uniformly at random. 
%Let $\mathcal R$ be the randomness space of $\A$, $a\getsr \A(b)$ is equivalent to $a=\A(b;r)$ for $r\getsr \mathcal R$.
By $\negl$ we denote an unspecified negligible function. %By $\poly$ we denote an unspecified polynomial function.


By $\vec x\in\bin^n$ we denote a column vector with size $n$,
% and by, say, $\vec x\in\one\times\bin^{n-1}$ we mean that the first element of $\vec x$ is $1$. 
 by $x_i$ we denote the $i$th element of a vector $\vec{x}$.
By $[n]$ we  denote the set $\{1, \cdots , n\}$, by $[t_1, t_2]$ we denote the set $\{t_1, t_1+1,\cdots ,t_2\}$, and by $\msb_k(\vec x)$ we denote the most significant $k$ bits in (the bit string representing) $\vec x$. By $\log n$ we mean $\lceil\log n\rceil$.
For some set or sequence $\set$, by $\vec x_\set$ we denote the sequence of all bits in $\vec x$ with indices (ordered lexicographically) in $\set$.
%By $||\vec x||$ we denote the hamming weight of $\vec x$. 
%By $\vec x_1\cdots,\vec x_\ell$ for some $\ell$, we denote $(\vec x_1^\top,\cdots,\vec x_\ell^\top)^\top$, i.e., the concatenation of $(\vec x_i)_{i\in[\ell]}$.
For a matrix $\matA$, 
we denote the set $\{\vec x \: | \: \matr A\vec x = \vec 0 \}$ by $\Ker(\matr A)$.  
%we denote the sets $\{\vec y \: | \: \exists \vec x \;\; \mbox{s.t.} \;\; \vec y = \matA\vec x \}$ and $\{\vec x \: | \: \matr A\vec x = \vec 0 \}$ by $\Ima(\matr A)$ (i.e., the span of $\matr A$) and $\Ker(\matr A)$ respectively.  
%By $\matr A^\bot\in\bin^{n\times (n-t')}$ we denote a non-zero matrix of rank $n-t'$ in $\Ker(\matr A^\top)$. Note that for any $\vec y\notin \Ima(\matr A)$, we have $\vec y^\top\matr A^\bot\neq \vec 0$. By $\overline{\matr A}$ (respectively, $\underline{\matr A}$) we denote the upper $(n-1)\times t$ matrix (respectively, lower $1\times t$ vector) of $\matr A$.
%Let $b\in \bin$, by $b\matr A$ we denote a zero matrix $\matr 0\in\bin^{n\times t}$ if $b=0$ or $\matr A$ if $b=1$.
%By $(a_{ij})_{i\in [l],j\in [m]}$ we denote the matrix 
%$\begin{pmatrix}a_{11} \cdots a_{1m}\\\vdots \ddots \vdots \\a_{l1}\cdots a_{lm}\\ \end{pmatrix}$. 
%Let $\matr A=(a_{ij})_{i\in [l],j\in [m]}$ be an $l\times m$ matrix and $\matr B=(\matr B_{ij})_{i\in [m],j\in [n]}$ be a large matrix consisting of $m\times n$ matrices $\matr B_{ij}$ for all $i\in[m]$ and $j\in[n]$. 
By $\matr I_n$ we denote an identity matrix in $\bin^{n\times n}$.
%By $\vec e_n^i$ we denote the column vector in $\bin^n$ with the $i$th element being $1$ and the other elements being $0$. 
By $\matr 0$ we denote a zero vector or matrix. 

By $H(x)$ we denote the \emph{Shannon entropy} of a random variable $x$ with alphabet $\mathcal{X}$ defined as
	$$H(x) = -\sum_{\hat x \in \mathcal{X}} \Pr[x=\hat x] \log \Pr(x=\hat x).$$
By $H(x\mid y)$ we denote the \emph{conditional entropy} of $x$ conditioned on a random variable $y$ with alphabet $\mathcal Y$ defined as
	$$
	H(x\mid y) = \sum_{\hat y \in \mathcal{Y}}\Pr[y=\hat y]\cdot H(x\mid y=\hat y).
	%&= - \mathbb{E}_{\hat x \in \mathcal{X}, \hat y \in \mathcal{Y}}\log\Pr(x=\hat x \big{|} y=\hat y).
	$$
By $I(x;y)$ we denote the \emph{mutual information} of $x$ and $y$ defined as $$I(x;y) = H(x) - H(x \mid y),$$ which is the reduction of the uncertainty of $x$ when $y$ is learned.
By $\expect[x]$ and $\var[x]$ we denote the \emph{expected value} and \emph{variance} of the variable $x$.
By $\tvd(x, y)$ we denote the \emph{total variation distance} of two random variables $x$, $y$ with alphabets $\mathcal{X}$ and $\mathcal{Y}$ defined as 
	$$ \tvd(x,y) = (1/2)\sum_{\hat{x} \in \mathcal{X} \cup \mathcal{Y}}|\Pr[x = \hat{x}]-\Pr[y = \hat{x}]|.$$
	A sequence of random variables $x_1,\cdots,x_q$ with alphabet $\mathcal X$ is said to be \emph{independent} if for any $\hat x_1,\cdots, \hat x_q \in \mathcal X$, we have
	$$\Pr[\wedge_{i=1}^q x_i = \hat x_i] = \prod_{i=1}^q\Pr[x_i = \hat x_i].$$
	 If we only require that 
	$$\Pr[x_i = \hat x_i\land x_j = \hat x_j] = \Pr[x_i = \hat x_i]\cdot \Pr[x_j = \hat x_j]$$
	 holds for any $1\leq i<j\leq q$, then the sequence is said to be \emph{pairwise independent}.
	
%\subsection{Probability Bounds}
%In this section, we introduce several probability bounds as follows.
%\begin{definition}[Pairwise independence]
%\label{def:pairwise}
%For any $q\in\N$, a sequence of random variables $x_1,\cdots,x_q$ with alphabet $\mathcal X$ are pairwise independent if for any $1\leq i<j\leq q$ and any $\hat x_1, \hat x_2 \in \mathcal X$, we have
%	$$\Pr[x_i = \hat x_1 \land x_j = \hat x_2] = \frac{1}{|\mathcal{Y}|^2}.$$
%\end{definition}
%
%\begin{definition}[Independence]
%\label{def:pairwise}
%For any $q\in\N$, a sequence of random variables $x_1,\cdots,x_q$ with alphabet $\mathcal X$ are pairwise independent if for any $\hat x_1, \cdots \hat x_q \in \mathcal X$, we have
%	$$\Pr[x_i = \hat x_1 \land x_j = \hat x_2] = \frac{1}{|\mathcal{Y}|^2}.$$
%\end{definition}

%\begin{theorem}[Chernoff bound]
%	Let $x = \Sigma_{i=1}^{n} x_i$, where $x_i = 1$ with probability $p_i$ and $x_i = 0$ with probability $1-p_i$, and all $x_i$ are independent. Let $\mu = \mathbb{E}[x] = \Sigma_{i=1}^{n}p_i$. Then 
%	\begin{compactitem}
%		\item $\Pr[x \geq (1+\delta)\mu] \leq e^{-\frac{\delta^2}{2+\delta}\mu}$ for all $\delta > 0$,
%		\item $\Pr[x \leq (1-\delta)\mu] \leq e^{-\frac{\delta^2}{2}\mu}$ for all $0 < \delta < 1$.
%	\end{compactitem}
%\end{theorem}





%\begin{definition}[Mutual information]
%	The \emph{mutual information} is defined as
%	$$I(x;y) = H(x) - H(x \big{|} y).$$
%\end{definition}
%By $\vec f_n^i$ we denote the vector in $\bin^{n}$ such that the first $i-1$ entries are $0$'s and the other entries are $1$'s. By $\matr E_n$ we denote the following $n\times (n-1)$ matrix, where the entries of the two main diagonals are $1$'s and the other entries are $0$'s.
%$$\matr E_n=
%\begin{pmatrix}
%1 &   &   &          &\\
%1 & 1 &   &          &\\
%%  & 1 & 1 &          &\\
%  &   &   & \ddots &\\
%  &   &   &          & 1 & 1\\
%  &   &   &          &   & 1
%\end{pmatrix}\in\bin^{n\times(n-1)}.
%$$
%One can check that ${\vec f_n^1}\in\Ker(\matr E_n^\top)$ and $\matr E_n\vec f_{n-1}^i=\vec e_n^i+\vec e_n^n$ for $i\in[n-1]$.%, and in $\matE\km^i\in\bin^{\lambda}$ for some $i\in[\lambda-1]$, the $i$th entry and the last entry are $1$'s and the other entries are $0$'s.
%By $\matr M^n_0$, $\matr M^n_1$, and $\matr N_n$, we denote the following $n \times n$ matrices:
%$\matr M^n_0 = \begin{pmatrix}\vec 0 &0\\ \matr I_{n-1} & \vec 0\end{pmatrix}$,
%$\matr M^n_1  = \begin{pmatrix}&\vec 0 &1\\ &\matr I_{n-1} & \vec 0\end{pmatrix}$,
%$\matr N_n=\begin{pmatrix}\vec 0& \matr 0\\1&\vec 0 \end{pmatrix}$.


%Let $X$ and $Y$ be random variables over a finite set $S$. The {\it statistical distance} between $X$ and $Y$ is defined to be
%\begin{equation*}
%\Dist(X,Y) = \frac{1}{2} \sum_{s \in S} |\Pr[X=s]-\Pr[Y=s]|.
%\end{equation*}
%We say that $X$ and $Y$ are $\epsilon$-close if $\Dist(X,Y) \le \epsilon$.




%For positive integers $k>1, \eta \in \Z^+$ and a matrix $\matr{A} \in \{0,1\}^{(k+\eta) \times k}$, we denote the upper square matrix of $\matr{A}$ by $\up{\matr{A}} \in \{0,1\}^{k \times k}$ and the lower $\eta$ rows of $\matr{A}$ by $\low{\matr{A}} \in \{0,1\}^{\eta \times k}$.
%We define $\Span(\matr A):=\{\matr A\vec r|\vec r\in\{0,1\}^k\}\subset \{0,1\}^{(k+\eta)}$.
%Similarly, for a column vector $\vec{v} \in \{0,1\}^{k+\eta}$, we denote the upper $k$ elements by  $\up{\vec{v}} \in \{0,1\}^{k}$ and the lower $\eta$ elements of $\vec{v}$ by $\low{\vec{v}} \in \{0,1\}^{\eta}$.
%For a bit string $\mes \in \bits^{n}$, $\mes_i$ denotes the $i$th bit of $\mes$ ($i\leq n$) and $\firstm i$ denotes the first $i$ bits of $\mes$.


%\heading{Games.} We follow~\cite{C:BlaKilPan14} to use code-based games for defining and proving security. A game $\gameg$ contains procedures $\initialize$ and $\finalize$, and some additional procedures $\procfont{P}_1,\ldots, \procfont{P}_n$, which are defined in pseudo-code. All variables in a game are initialized as 0, and all sets are empty (denote by $\emptyset$).
%An adversary $\advA=\{\adv\}_{\lambda\in\N}$ is executed in game $\gameg$ w.r.t. the security parameter $\lambda$ (denote by ${\gameg}^{\adv}$) if $\adv$ first calls $\initialize$, obtaining its output. Next, it may make arbitrary queries to $\procfont{P}_i$ (according to their specification) and obtain their output. Finally, it makes one single call to $\finalize(\cdot)$ and stops. We use $ \gameg^{\adv} \outp d$ to denote that $\gameg$ outputs $d$ after interacting with ${\adv}$, and $d$ is the output of $\finalize$.




%

%% entropy definitions

%end entropy definitions


%\subsection{Function Families and Word-Random Access Machine}
%
%%We now recall the definitions of function family. .
%
%\begin{definition}[Function family]
%
%\end{definition}

%\begin{definition}[$\nczero$]
%\label{def:nczero}
%The class of (non-uniform) $\aczero$ function families is the set of all function families $\mathcal{F} = \{f_{\lambda}\}_{\lambda\in\N}$ for which there is a polynomial $p(\cdot)$ and constant $d$ such that for each $\lambda$, $f_{\lambda}$ can be computed by a (randomized) circuit of size $p(\lambda)$, depth $d$, and fan-in 2 using {$\AND$}, {$\OR$}, and {$\NOT$} gates.
%\end{definition}

%\begin{definition}[$\aczero$]
%\label{def:aczero}
%The class of (non-uniform) $\aczero$ function families is the set of all function families $\mathcal{F} = \{f_{\lambda}\}_{\lambda\in\N}$ for which there is a polynomial $p(\cdot)$ and constant $d$ such that for each $\lambda$, $f_{\lambda}$ can be computed by a (randomized) circuit of size $p(\lambda)$, depth $d$, and unbounded fan-in using {$\AND$}, {$\OR$}, and {$\NOT$} gates.
%\end{definition}
%One can easily see that $\nczero$ is a subset of $\aczero$, and for any polynomial $n=n(\lambda)$ and $\vec x,\vec y\in\bin^n$ where either $\vec x$ or $\vec y$ has only constant Hamming weight, the inner product of $\vec x$ and $\vec y$ is computable in $\nczero$.

%Notice that we can assume that an $\aczero$ circuit has no $\NOT$ gates since we can remove all $\NOT$ gates to just the input, according to the De Morgan Rules. Moreover, we have $\nczero\subsetneq\aczero$ by known results.

%\begin{definition}[Sparse vector]
%Let $\lambda$ be the security parameter. A vector $\vec x$ is said to be \emph{sparse} if there exists a constant $c$ and a function $t(\cdot)$ such that $t(\lambda)\leq O(\log^c\lambda)$ and $||\vec x||=t(\lambda)$.
%\end{definition}

%\begin{lemma}[Polylog inner products \cite{C:DegVaiVas16}]
%\label{lem:polylog-ip}
%For any constant $c$ and any function $t(\cdot)$ such that $t(\cdot)\leq O(\log^c\lambda)$, there exists a function family $\{\ip:\bin^\lambda\times\bin^\lambda\}_{\lambda\in\N}$ such that for any $\vec x,\vec y\in\bin^\lambda$, where $\min(||x||,||y||)\leq t(\lambda)$, $\ip(x,y)=\langle x,y \rangle$.
%\end{lemma}

%\begin{lemma}[Polylog inner products \cite{C:DegVaiVas16}]
%\label{lem:polylog-ip}
%There exists a function family $\{\ip:\bin^\lambda\times\bin^\lambda\}_{\lambda\in\N}\in\aczero$ such that 
%for any $\vec x,\vec y\in\bin^\lambda$, where at least one of $\vec x$ and $\vec y$ is sparse, we have $\ip(x,y)=\langle x,y \rangle$.\end{lemma}
%Notice that for any polynomial $n = n(\lambda)$, we have $\{f_n\}_{\lambda\in\N} \in \ncone$ iff $\{f_\lambda\}_{\lambda\in\N}\in\ncone$. Hence, we can also replace $\lambda$ by $n$ in the above lemma.

%Let $\{\PARITY\}_{\lambda\in\N}$ be the function family such that for all $\lambda\in\N$, $\PARITY$ on input any $\vec x\in\bin^\lambda$ outputs $\sum_{i=1}^\lambda x_i$. The following theorem states that any $\aczero$ circuit has very small correlation with $\PARITY$.
%\begin{theorem}[\cite{Has14,IMP12}]
%\label{thm:correlation}
%For any $\A = \{a_{\lambda}\}_{\lambda\in\N}\in\aczero$ with size $p$ and constant depth $d$ and any $\lambda\in\N$, we have
%\[\begin{split}
%\Bigg|\Pr_{\vec x\getsr\bin^\lambda}&[a_{\lambda}(\vec x)=1|\PARITY(\vec x)=1]\\
%&-\Pr_\vec {x\getsr\bin^\lambda}[a_{\lambda}(\vec x)=1|\PARITY(\vec x)=0] \Bigg|\leq 2^{-\Omega(\lambda/\log^{d-1}(p))}.\end{split}
%\]
%\end{theorem}
%One can see that for any polynomial $p$ in $\lambda$, $2^{-\Omega(\lambda/\log^{d-1}(p))}=2^{-\Omega(\lambda/\log^{d-1}(\lambda))}$ is negligible.




%\begin{proof}
%%First, we note that $\matE\vec r$ for $\vec r\getsr\bin^\lambda$ is uniformly distributed in $\bin^\lambda$ subject to the constraint that $\PARITY(\matE\vec r)=0$.
%Let $\A=\{\adv\}_\lambda\in\aczero$ be any adversary that distinguishing $\matE\vec r$ and $\matE\vec r+\vece$ for $\vec r\getsr\bin^\lambda$ with non-negligible advantage, we show $\{\PARITY\}_{\lambda\in\N}\in\aczero$, which leads to a contradiction.
%
%First, we note that for any vector $\vec x\in\bin^\lambda$, the following distributions are identical.
%$$\matE\vec r+\vec x\ \mbox{and}\ \matE\vec r+\vece \cdot b,$$
%where $\vec r\getsr\bin^\lambda$ and $b=\PARITY(\vec x)$. This is due to the fact that the first $\lambda-1$ bits of a vector $\vec y=(y_i)_{i=1}^\lambda$ sampled from either distribution are uniformly distributed (since $\overline{\matr E}_\lambda$ is of full rank), and $y_\lambda=\sum_{i=1}^{\lambda-1}y_i$ must hold in both cases.
%Hence, we have
%$$|\Pr[\adv(\matE\vec r+\vec x)=1|\PARITY(\vec x)=1]=\Pr[a_{\lambda}(\matE\vec r+\vece)=1]$$
%and 
%$$|\Pr[\adv(\matE\vec r+\vec x)=1|\PARITY(\vec x)=0]=\Pr[a_{\lambda}(\matE\vec r)=1]$$
%for $\vec r\getsr\bin^\lambda$. Let the former probability be $p$ and the latter be $q$. To decide $\PARITY(\vec x)$ given any $\vec x\in\bin^\lambda$, we just have to determine whether $\Pr[\adv(\matE\vec r+\vec x)|\vec r\getsr\bin^{\lambda-1}]<\frac{p+q}{2}$ by taking samples from $\adv(\matE\vec r+\vec x)$ in parallel and using the threshold function to check whether more than a $\frac{p+q}{2}$ fraction of the samples are $1$. Due to the fact $p$ and $q$ are non-negligibly separated, $\PARITY(\vec x)$ can be decided from polynomial number of samples with exponentially small failure probability. Also, all the operations can be performed in $\aczero$ since all the row vectors in $\matE$ have only constant hamming weight, i.e., $\matE\vec r$ can be computed in $\aczero$, completing the proof.
%\end{proof}


%\begin{definition}[$\maczero$]
%The class of (non-uniform) $\maczero$ function families is the set of all function families $\mathcal{F} = \{f_{\lambda}\}_{\lambda\in\N}$ for which there is a polynomial $p(\cdot)$ and constant $d$ such that for each $\lambda$, $f_{\lambda}$ can be computed by a (randomized) circuit of size $p(\lambda)$, depth $d$, and unbounded fan-in using {$\AND$} and {$\OR$} gates.
%\end{definition}







\subsection{Fine-Grained Multi-Party NIKE}
\label{def:munike}
In this section, we define fine-grained multi-party NIKE and multi-party key exchange in the BSM. In the following, 
a circuit (respectively, word-RAM) family is a family of (possibly randomized) circuits (respectively, word-RAMs) with respect to all security parameters $\lambda\in\N$, denoted by $\{f_{\lambda}\}_{\lambda\in\N}$. For brevity, we refer to it as $\{f_{\lambda}\}$.
%In this section, we give the definitions of $\Sigma$-protocol, NIZK, and non-interactive zap. 
%Below, for a language description $\rho$ with the associated language $\lang_\rho$ and relation $\relation_\rho$, by $\x\in\lang_\rho$ we mean that there exists $\w$ such that $\relation_\rho(\x,\w)=1$.
%\heading{NIZK.} We now give the definition of fine-grained NIZK with composable zero-knowledge and statistical/perfect soundness.
%Let $\dist_{\lambda}$ be a probability distribution over a collection of relations $\Relation=\{\relation_{\matr M}\}_{\matr M\in\dist_{\lambda}}$ parametrized by a matrix $\matr M\in \bin^{n\times t}$ of rank $t'\leq n$ generated as $(\matr M^\top,\matr M^\bot) \getsr \dist_{\lambda}$ with the associated language 
%\[{\mathcal L}_{\matr M}=\{ {\vec t}:\exists \vec w \in \bin^t, \text{ s.t. } \vec t= \matr M \vec w\}.\]
%Notice that similar to witness sampleable distribution in the classical world~\cite{AC:JutRoy13}, we require that $\dist_\lambda$ additionally outputs a non-zero matrix $\matr M^\bot\in\bin^{n\times (n-t')}$ in the kernel of $\matr M^\top$. %We note that the matrix distribution in \cref{def:matrdist} is sampleable. 
%An example of sampleable distribution is $\ZeroSamp(n)$, which can additionally sample a non-zero vector in the kernel of its output. \footnote{In fact, the rightmost vector $(r_1,\cdots,r_{n-1},1)^\top$ of the intermediate matrix generated by $\RSamp(n)$ (see \cref{fig:lsamp}) forms a vector in the kernel of $\matr M^\top$. See the proof of Lemma 3 in \cite{AC:EgaWanTan19} for more details}

%We define the notions of $\QANIZK$, designated-prover $\QANIZK$ ($\DPQANIZK$), designated-verifier $\QANIZK$ (\( \DVQANIZK \)), designated-prover-verifier $\QANIZK$ (\( \DPVQANIZK \)) as follow.

We now define the multi-party NIKE and its security. We require that the users be in $\mathcal C_1$ and the adversaries be in $\mathcal C_2$. In this work, $\mathcal C_1$ and $\mathcal C_2$ represent classes of circuit families in the bounded-parallel-time model, and represent classes of word-RAM families in the bounded-time model. In the BSM, $\mathcal C_1$ refers to a class of streaming word-RAM families with bounded storage, which receive inputs in a streaming manner and store the outputs.
\begin{definition}[Multi-party non-interactive key exchange]
\label{def:nike}
A \emph{$\cc_1$-$\pnum$-party non-interactive key exchange scheme (NIKE)} consists of two algorithm families $\mathsf{NIKE}=\{\gen,\allowbreak\share\}\in\cc_1$ such that
	\begin{compactitem}
	\item $\gen$ on input $\urs\in\bin^n$ for some $n=n(\lambda)$ outputs a public key $\pk$ and a secret key $\sk$. 
	\item $\share$ is a deterministic algorithm that on input a set of public keys $(\pk_i)_{i\in\cs}$, an index $i$, and a secret key $\sk_i$ outputs either $\key\in\mathcal K_\lambda$ or the abort symbol $\bot\notin\mathcal K_\lambda$, where $i\in\cs$ and $\mathcal K_\lambda$ is the session key space.
	\end{compactitem}
	
	\emph{$\epsilon$-correctness} is satisfied if for all $\lambda\in\N$, all $i,j\in\cs$, we have
	$$\Pr[\exists \key\in\mathcal K_\lambda:\share(i,\sk_i,(\pk_i)_{i\in\cs})=\share(j,\sk_j,(\pk_i)_{i\in\cs})=\key]\geq 1-\epsilon
	$$
	where
	$\urs\getsr\bin^n$ and $(\pk_i,\sk_i)\getsr\Gen(\urs)$ for all $i\in\cs$.
%		 is satisfied if for any adversary $\advA=\{\adv\}\in\cc_2$ and an infinite number of values $\lambda\in\N$, we have
%$$\Pr[\key\getsr \adv((\pk_i)_{i\in\cs})]\leq\epsilon,$$
%where $(\pk_i,\sk_i)\getsr\Gen$ for all $i\in\cs$, $\key=\share(j,\sk_j,(\pk_i)_{i\in\cs})$ for some $j\in[K]$, and $\key'\getsr\mathcal K_\lambda$.

	\emph{$\cc_2$-$\epsilon$-security} (respectively, \emph{$\cc_2$-$\epsilon$-weak security}) is satisfied if for any adversary $\advA=\{\adv\}\in\cc_2$ and all sufficiently large $\lambda\in\N$, we have
$$|\Pr[1\getsr \adv(\urs,(\pk_i)_{i\in\cs},\key)]-\Pr[1\getsr \adv(\urs,(\pk_i)_{i\in\cs},\key')]|\leq\epsilon$$
$$(\mbox{respectively,}\ \Pr[\key\getsr \adv(\urs,(\pk_i)_{i\in\cs})]\leq\epsilon)$$
where $\urs\getsr\bin^n$, $(\pk_i,\sk_i)\getsr\Gen(\urs)$ for all $i\in\cs$, $\key=\share(1,\sk_1,(\pk_i)_{i\in\cs})$, and $\key'\getsr\mathcal K_\lambda$.
In the plain model, we use the convention $\urs=\bot$.
If $\share$ outputs $\bot$ in a security experiment, the actual value $\bot$ is given to the adversary; thus abort probabilities are accounted for by the corresponding correctness and security bounds.

\emph{$(\amem,\epsilon,\theta)$-security in the BSM} is satisfied if for any function $f:\bin^n\rightarrow\bin^\amem$, we have
$$\exists \event:\ \Pr[\event]\geq 1-\epsilon\ \mbox{and}\ I(\key;(f(\urs),(\pk_i)_{i\in[\pnum]})\mid\event)\leq \theta,$$
where $\urs\getsr\bin^n$, $(\pk_i,\sk_i)\getsr\Gen(\urs)$ for all $i\in\cs$, and $\key=\share(1,\sk_1,\allowbreak(\pk_i)_{i\in\cs})$.

%\begin{definition}[$(\mu,\mu')$-local key security] 
%\end{definition}

\end{definition}

\heading{Remarks on the BSM.} 
Here we note that the security in the BSM captures the situation where
all parties have access to $\urs$ and the adversary is allowed to apply any (potentially unbounded) storage function $f$ with an output length limit of $\amem$ and store information on $\urs$. Afterwards, $\urs$ becomes inaccessible, and all parties broadcast public keys, which are typically required to be short, and share a session key in a non-interactive way.
Notice that we consider the traditional BSM \cite{C:CacMau97} rather than the streaming BSM proposed by Dodis, Quach, and Wichs~\cite{EC:DodQuaWic23}. In the streaming BSM, the parties are allowed to communicate interactively with perhaps multiple rounds of long messages.

\heading{Extendability of Multi-Party NIKE.}
%Roughly, The extendability of NIKE is satisfied if the sharing algorithm of NIKE can be decomposed into two algorithms. The first algorithm, given a private key and partial public keys as input, outputs a partial private key. The second algorithm, taking the partial private key and a public key as input, generates an output almost identical to the output of the original sharing algorithm.
%Roughly, the extendability of a multi-party NIKE says that the sharing procedure involves combining the secret key with public keys of other users sequentially. Each intermediate value, termed as a “local key”, essentially represents a shared key (without privacy amplification) amongst a subgroup of all parties. 
We now define a property of multi-party NIKE called extendability.
Roughly, the extendability of a multi-party NIKE says that the sharing procedure involves an algorithm termed as $\combine$, which sequentially combines a user's secret key with other public keys, and the intermediate value produced by $\combine$, referred to as a local key, can be used to derive the shared key by using another algorithm $\extract$. %This process results in a shared key for a multi-party NIKE if $\combine$ combines a secret key using the public keys of all other users.
\begin{definition}[Extendability]
\label{def:ext}
	A $\cc_1$-$\pnum$-party non-interactive key exchange scheme $\mathsf{NIKE}=\{\gen,\allowbreak\share\} \in \cc_1$ satisfies \emph{$\epsilon$-extendability} if there exist two deterministic algorithm families $\{\combine, \allowbreak\extract \} \in \cc_1$ such that, for every $\lambda\in\N$, the following holds. Let
	\[
	\urs\getsr\bin^n\ \mbox{and}\
	(\pk_i,\sk_i)\getsr\gen(\urs)\quad\text{for all}\ i\in[\pnum].
	\]
	There exists an event $\event$ with $\Pr[\event]\geq 1-\epsilon$ such that, conditioned on $\event$, for every $\indexset \subseteq [\pnum]$, nonempty $\indexset' \subseteq \indexset$, $j \in \indexset$, and $i^{\prime} \in \indexset'$, we have
%	\item for any set $\indexset \subseteq [\pnum-1]$ and $i^{\prime}, j^{\prime} \in \indexset$, we have
%	$$ \combine^1(\indexset, \sk_{i^{\prime}}, (\pk_i)_{i \in \indexset}) = \combine^1(\indexset, \sk_{j^{\prime}}, (\pk_i)_{i \in \indexset}),$$
%	where $(\pk_i, \sk_i) \getsr \gen$ for any $i \in [\pnum]$.
%<<<<<<< HEAD
%	\[\begin{split}\combine&(\indexset, \indexset',\localkey_{\indexset'}, \{\pk_i\}_{i \in \indexset\slash\indexset'}) =
%	 \combine(\indexset, \{j\}, \sk_j, \{\pk_i\}_{i \in \indexset\slash\{j\}})\end{split}\]
	\[\begin{split}\combine&(\localkey_{\indexset'}, \{\pk_i\}_{i \in \indexset\backslash\indexset'}) =
	 \combine(\sk_j, \{\pk_i\}_{i \in \indexset\backslash\{j\}})\end{split}\]
	%\item%=======
%	\item for any set $\indexset \subseteq [\pnum]$, $\indexset_1 \subseteq \indexset$, and $j \in \indexset$, let $\indexset_2 = \{j\}$, $\localkey_2 = \sk_{j}$, and $\localkey_1 = \combine(\indexset_1, \{i^{\prime}\}, \sk_{i^{\prime}}, \allowbreak \{\pk_i\}_{i \allowbreak \in \indexset_1\slash\{i^{\prime}\}})$ for any $i^{\prime} \in \indexset_1$, we have
%	$$ \combine(\indexset, \indexset_1, \localkey_1, \{\pk_i\}_{i \in \indexset\slash\indexset_1}) = \combine(\indexset, \indexset_2, \localkey_2, \{\pk_i\}_{i \in \indexset\slash\indexset_2}),$$
%	\item let $\localkey = \combine([\pnum], \{j\}, \sk_{j}, \{\pk_i\}_{i \in [\pnum]\slash\{j\}})$ for any $j \in [\pnum]$ we have
%>>>>>>> refs/remotes/origin/main
	and
	\[
	\extract(\localkey_{[\pnum]}, \pk_{\pnum}) = \share(1,\sk_1,(\pk_i)_{i\in\cs}),
	\]
	%\end{compactitem}
	where $\localkey_{\indexset'} = \combine(\sk_{i^{\prime}}, \allowbreak \{\pk_i\}_{i \allowbreak \in \indexset'\backslash\{i^{\prime}\}})$ and $\localkey_{[\pnum]} = \combine(\sk_{1}, \{\pk_i\}_{i >1})$.
\end{definition}
%We also give the definition of restricted extendability, which is sufficient for constructing the restricted BPRF.
We also give the definition of restricted extendability, which requires that $\combine$ only generates the local key of the users with consecutive indices.

%\heading{Restricted extendability.} For a $\cc_1$-$\pnum$-party non-interactive key exchange scheme $\mathsf{NIKE}=\{\gen,\allowbreak\share\} \in \cc_1$, the restricted $\epsilon$-extendability is the special case of $\epsilon$-extendability where we require the set $\indexset = [\leftindex+1, \pnum-1-\rightindex] \subseteq [\pnum-1]$ for some $\leftindex, \rightindex \in [0, \pnum-1]$ such that $\leftindex + \rightindex \leq \pnum-1$.
%\begin{definition}[Restricted extendability] 
%\label{def:resext}
%\emph{$(\epsilon,\epsilon')$-restricted extendability} is defined in exactly the same way as $(\epsilon,\epsilon')$-extendability except that the set $\indexset \subseteq [\pnum]$ is restricted to $\indexset = [\leftindex, \rightindex] \subseteq [\pnum]$ for any $\leftindex, \rightindex \in [\pnum]$ and $\indexset'$ is restricted to $\indexset' = [\subleftsize, \subrightsize]$ for any $\subleftsize, \subrightsize \in [\pnum]$ such that $\rightsize\geq\subrightsize\geq\subleftsize \geq \leftsize$.
%\end{definition}
\begin{definition}[Restricted extendability] 
\label{def:resext}
\emph{$\epsilon$-restricted extendability} is defined in exactly the same way as $\epsilon$-extendability except that the set $\indexset \subseteq [\pnum]$ is restricted to $\indexset = [\leftindex, \rightindex] \subseteq [\pnum]$ for any $\leftindex, \rightindex \in [\pnum]$ and $\indexset'$ is restricted to $\indexset' = [\subleftsize, \subrightsize]$ for any $\subleftsize, \subrightsize \in [\pnum]$ such that $\rightsize\geq\subrightsize\geq\subleftsize \geq \leftsize$.
\end{definition}


\heading{Local key independence in the BSM.} Next we define a natural property called local key independence for NIKE in the BSM.
Roughly, it says that public information and auxiliary key pairs reveal only limited
information about local keys, and that adding one more independently generated key
pair reveals only limited additional information through the common URS. %\begin{definition}[Local key independence]
%A $\pnum$-party NIKE in the BSM with (restricted) extendability satisfies
%\emph{$\mu$-local key independence in the BSM} if we have
%%$$\exists \event:\ \Pr[\event]\geq 1-\delta\ \mbox{and}\ I(\localkey_{[i]};(\pk_j,\pk_j',\sk_j')_{j\in[n]})\leq \mu,$$
%$$I(\localkey_{[i]};(\pk_j,\pk_j',\sk_j')_{j\in[\pnum]})\leq \mu,$$
%where $\urs\getsr\bin^n$, $(\pk_j,\sk_j)\getsr\Gen(\urs)$, $(\pk_j',\sk_j')\getsr\Gen(\urs)$ for all $j\in\cs$, and $\localkey_{[i]}=\combine(\sk_1,(\pk_j)_{2\leq j\leq i\}})$.
%\end{definition}
\begin{definition}[Local key independence]
\label{def:localkeyind}
A $\pnum$-party NIKE in the BSM with $\epsilon$-(restricted) extendability satisfies
\emph{$((\mu_i)_{i\in[\pnum]},t)$-local key independence in the BSM}, for $t\in\N$, if the following holds. In the key-generation experiment of \cref{def:ext} (or \cref{def:resext}), let $\event$ be the corresponding extendability event. Conditioned on the same $\urs$, independently sample $(\pk_j',\sk_j')\getsr\Gen(\urs)$ for $j\in[t]$ and set $\localkey_{[i]}=\combine(\sk_1,(\pk_j)_{2\leq j\leq i})$. Then, for every $i\in[\pnum]$, we have
\[
I(\localkey_{[i]};(\pk_j)_{j\in[i]},(\pk_j',\sk_j')_{j \in [t]}\mid \event)\leq \mu_i,
\]
and for every $i\in[\pnum]$ with $i>1$, we have
\[
\begin{split}
&I(\localkey_{[i-1]};(\pk_i,\sk_i)\mid(\pk_j)_{j\in[i-1]},(\pk_j',\sk_j')_{j\in[t]},\event)\\
&\leq
I(\localkey_{[i-1]};(\pk_i,\sk_i)\mid(\pk_j)_{j\in[i-1]},(\pk_j',\sk_j')_{j\in[t]},\urs,\event)+\mu_i.
\end{split}
\]
\end{definition}
%The second condition only bounds the residual dependence caused by the common URS. After conditioning on the URS, the remaining averaging step is handled by \cref{lem:DM04}.

%Also, depending on the specific applications, it might be anticipated that one of the parties will use only a small amount of memory, denoted as $\umem'$.

%\begin{definition}[Multi-party key agreement in the BSM]
%\label{def:bsmnike}
%A \emph{$((\umem,\umem'),\ucom)$-$\pnum$-party key agreement} in the bounded storage model (BSM) is a protocol between $\pnum$ parties $(\party_i)_{i\in[\pnum]}$ where the memory of one party is bounded by $\umem'$ and the memory of each of the remaining parties is bounded by $\umem$.
%	\begin{compactitem}
%	\item For each $i\in[\pnum]$, $\party_i$ observes a random string $\urs\getsr\bin^n$ for some $n$ in a streaming manner, and generates a key pair $(\pk_i,\sk_i)$ and stores it. 
%	\item When $\urs$ disappeared, each party $(\party_i)_{i\in[\pnum]}$ publishes $\pk_i$, and each $\party_i$ generates a session key $\key_i\in\bin^k$.
%	\end{compactitem}
%	
%	\emph{$\epsilon_c$-correctness} is satisfied if we have
%	$\Pr[\key_1=\cdots=\key_\pnum]\geq 1-\epsilon_c.
%	$
%	and the communication complexity between users is bounded by $\ucom$, i.e., $|\trans|\leq \ucom$
%	when the above protocol is executed honestly.
%	
%\emph{$(\amem,\epsilon,\theta)$-security} is satisfied if for any storage function $\advA$ with domain $\bin^n$ and range $\bin^\amem$, we have
%$$\exists \event:\ \Pr[\event]\geq 1-\epsilon\ \mbox{and}\ I(\key;(\A(\urs),\trans,\event))\leq \theta.$$
%\end{definition}


%\begin{definition}[Adaptive security]
%\label{def:as}
%$\nike$ is said to satisfy \emph{adaptive security} if for any adversary $\advA=\{\adv\}$, an infinite number of values $\lambda\in\N$, and any polynomial number $n=n(\lambda)$, we have
%$$|\Pr[\AS_\real^{\adv,n}\Rightarrow 1]-\Pr[\AS_\rand^{\adv,n}\Rightarrow 1]|\leq\negl(\lambda),$$
%where the experiments are defined as in \cref{fig:sec-as}.
%
%\begin{figure}[h]
%\begin{center}
%\framebox{
%\small
%  \begin{tabular}{l|l}
%    \begin{minipage}[t]{0.45\textwidth}
%    
%                    \underline{$\init$:}\\
%    $(\pk_i,\sk_i)\getsr\gen$ for all $i\in[n]$\\
%    
%    \underline{$\ShareO(\cs)$:}\\
%    $\Qc\leftarrow \Qc\cup\{\cs\}$\\
%    $\key_\cs=\share(i,\sk_i,(\pk_i)_{i\in\cs})$ for some $i\in\cs$\\
%    Return $\key_\cs$
%	
%	\end{minipage} & \begin{minipage}[t]{0.45\textwidth}
%
%	
%	\underline{$\ShareO^*(\cs^*)$:}\\
%    $\key=\share(\pk_{\cs^*},i,\sk_i)$ for some $i\in\cs^*$\\
%    \fbox{$\key\getsr\mathcal K$}\\
%	
%       \underline{$\finalize(\beta)$:}\\
%       If $\cs^*\notin\Qc$, return $\beta$\\
%       Else return $0$
%    \end{minipage} 
%      \end{tabular}
%      }
%\end{center}
%\caption{Definition of $\AS_\real$ and \fbox{$\AS_\rand$}.}
%\label{fig:sec-as}
%\end{figure}
%\end{definition}
%
%\begin{definition}[Strengthened adaptive security]
%\label{def:sas}
%$\nike$ is said to satisfy \emph{strengthened adaptive security} if for any adversary $\advA=\{\adv\}$, an infinite number of values $\lambda\in\N$, and any polynomial number $n=n(\lambda)$, we have
%$$|\Pr[\SAS_\real^{\adv,n}\Rightarrow 1]-\Pr[\SAS_\rand^{\adv,n}\Rightarrow 1]|\leq\negl(\lambda),$$
%where the experiments are defined as in \cref{fig:sec-sas}.
%
%\begin{figure}[h]
%\begin{center}
%\framebox{
%\small
%  \begin{tabular}{l|l}
%    \begin{minipage}[t]{0.4\textwidth}
%    
%                    \underline{$\init$:}\\
%    $(\pk_i,\sk_i)\getsr\gen$ for all $i\in[n]$\\
%    
%    \underline{$\share(i)$:}\\
%    $\Qc\leftarrow \Qc\cup\{i\}$\\
%    Return $\key_i$
%	
%	\end{minipage} & \begin{minipage}[t]{0.5\textwidth}
%
%	
%	\underline{$\ShareO^*(\cs)$:}\\
%    $\key=\share((\pk_i)_{i\in\cs},i,\sk_i)$ for some $i\in\cs$\\
%    \fbox{$\key\getsr\mathcal K$}\\
%	
%       \underline{$\finalize(\beta)$:}\\
%       If $i\notin\cs$ for all $i\in\Qc$, return $\beta$\\
%       Else return $0$
%    \end{minipage} 
%      \end{tabular}
%      }
%\end{center}
%\caption{Definition of $\SAS_\real$ and \fbox{$\SAS_\rand$}.}
%\label{fig:sec-sas}
%\end{figure}
%\end{definition}
